\section{Thiết lập thực nghiệm}
\subsection{Train và chọn checkpoint}
Các lần train dùng seed 42, 43 và 44 trên cùng manifest. Mỗi classifier chạy 20 epoch, batch size 32, optimizer AdamW, learning rate $10^{-3}$, weight decay $10^{-4}$ và cross-entropy loss. Nhánh Forensic dùng batch ảnh 64 khi trích xuất FFT.

Checkpoint được chọn theo AUROC validation. Nếu AUROC bằng nhau, BAcc tại ngưỡng 0.5 được dùng để chọn. Sau khi chọn checkpoint, ngưỡng được chọn theo BAcc validation cao nhất, tiếp đến F1 cao nhất và gần 0.5 nhất nếu còn bằng nhau. Fusion cũng chọn ngưỡng trên validation. Các lần thử ảnh sau xử lý giữ nguyên checkpoint và ngưỡng này.

Cấu hình cuối được xếp hạng theo macro AUROC validation trung bình qua ba seed, rồi BAcc trung bình, AUROC SD nhỏ hơn và AUROC của generator yếu nhất. Quy tắc này chọn cấu hình, không lấy seed tốt nhất làm kết quả chính hoặc ensemble ba seed.

\subsection{Cách tính điểm}
Dự đoán nhãn là $\hat y=\mathbf{1}[p(x)\geq t]$. AUROC được tính từ xác suất liên tục. Balanced Accuracy được tính bằng:
\begin{equation}
\mathrm{BAcc}=\frac12\left(\frac{\mathrm{TP}}{\mathrm{TP}+\mathrm{FN}}+
\frac{\mathrm{TN}}{\mathrm{TN}+\mathrm{FP}}\right).
\end{equation}
Accuracy và F1 được tính thêm để so sánh. Mỗi dòng theo generator gồm 2.000 ảnh AI và cùng 4.000 ảnh thật, tổng 6.000 ảnh. Điểm macro là trung bình của hai dòng generator, còn điểm overall dùng 8.000 ảnh test duy nhất.

Do hai generator có cùng số ảnh AI và dùng chung ảnh thật, macro AUROC và BAcc bằng điểm overall trong tập test này. Accuracy và F1 không có tính chất đó.

Kết quả được báo cáo bằng mean và sample SD qua ba seed, với $ddof=1$. Ba seed dùng cùng dữ liệu, nên SD mô tả chênh lệch giữa các lần chạy trên tập pilot.

\subsection{Thử ảnh sau xử lý}
Đánh giá gồm ảnh sạch và bốn phép biến đổi:
\begin{itemize}[nosep]
\item Nén JPEG với quality 70 và subsampling 2.
\item Resize về $112\times112$ rồi tăng lại $224\times224$ bằng bilinear.
\item Gaussian blur với radius 1.0.
\item Crop vuông ở tâm, cạnh bằng $\lfloor0.8\min(W,H)\rfloor$, rồi resize về $224\times224$.
\end{itemize}
Mỗi biến đổi được áp dụng riêng từ ảnh RGB gốc. Crop 80\% tính theo cạnh ngắn của ảnh. Thí nghiệm chưa thử nhiều mức nén, blur hoặc crop. Mỗi seed có $5\times6=30$ kết quả theo điều kiện và mô hình, tổng cộng 90.

\subsection{Bootstrap và ảnh fusion sửa đúng hoặc làm sai}
Khoảng tin cậy được tính bằng percentile bootstrap 95\% với 2.000 lần lấy mẫu, seed 2026. Prediction được lấy mẫu lại theo nhóm nguồn ảnh và nhãn. Các mô hình trong cùng phép so sánh dùng chung chỉ số lấy mẫu, còn 4.000 ảnh thật vẫn được dùng chung khi tính điểm hai generator.

Ba cặp đã kiểm tra là L+F so với L, G+L+F so với G+L và F so với L, cho AUROC và BAcc ở cả năm điều kiện. Kết quả gồm 180 dòng khoảng tin cậy theo mô hình và điều kiện, cùng 90 dòng cho chênh lệch giữa các cặp. Bootstrap chỉ dùng prediction đã lưu, không train hoặc suy luận lại.

Rescue là ảnh baseline đoán sai nhưng fusion đoán đúng. Harm là trường hợp ngược lại. Net gain bằng rescue trừ harm, cho biết số ảnh đúng thêm sau fusion. Mỗi cấu hình dùng ngưỡng riêng đã chọn trên validation.

\subsection{Kiểm tra đặc trưng đơn giản và môi trường chạy}
Logistic regression được train ở seed 42 để xem đặc trưng đơn giản có phân biệt được hai lớp không. Một cấu hình chỉ dùng dung lượng file. Cấu hình còn lại thêm mean và standard deviation độ sáng, độ bão hòa, tỷ lệ pixel xám và đặc trưng cạnh. Model được train trên train và đánh giá trên validation.

Metadata seed 42 ghi baseline chạy CUDA, Forensic chạy Tesla T4 và fusion chạy CPU. File môi trường Forensic ghi Python 3.13.15 và PyTorch 2.11.0+cu128. Fusion dùng Python 3.13.15 và PyTorch 2.11.0+cpu. Báo cáo Phase F ghi backend suy luận MPS cho seed 42/43 và CUDA cho seed 44. Đây là thông tin backend của từng giai đoạn, không phải xác nhận thiết bị train cho tất cả các seed.
